\section*{Expected Outcome}

My main hypothesis (H1) is that the change in capital expenditure associated with exposure to Helene is more positive for firms with higher cash holdings before the storm. Cash can finance replacement assets and help maintain planned investment when operating receipts fall. It can also cover the period before a firm obtains external finance or receives insurance payments.

Firms with more cash may increase capital expenditure more or reduce it less, so H1 does not require an overall rise in investment spending after Helene. This comparison is most informative when firms face similar investment needs. Insurance and access to credit may reduce the importance of cash, while shortages of labour or equipment may delay investment even when funds are available.
